% Study sheet: Volume I, Chapter 15, Responsible Engineering.
% Numbers: \val{...} from responsible_engr_studysheet_numbers.tex, written by responsible_engr_studysheet_numbers.py.
\documentclass{studysheet}
\input{responsible_engr_studysheet_numbers}
\sheetvolume{1}
\sheetchapter{15}{Responsible Engineering}
\sheetsources{Numbers from the chapter's constructed credit service and DLRM ranking scenarios (mlsysim); assumed inputs are marked.}

\begin{document}
% ===================================================================== FRONT: LEARN
\frontside
\begin{bigidea}
A model can meet every latency and accuracy target and still harm people, because its objective is an
\hl{incomplete specification}. Responsible engineering writes the missing requirements as \hl{measurable bounds}, tests them before release, and keeps
measuring them after.
\end{bigidea}

\begin{leftcol}
\sheetsec{Key ideas}
\begin{ideas}
\idea{Verified is not validated.}{A model can fit its objective exactly while that proxy drifts
  from the goal. \emph{Alignment gap}: proxy minus goal on a common baseline.}
\idea{Locate the failure on D·A·M.}{Data: biased records or drift on a slice. Algorithm: a misaligned
  target or loss. Machine: telemetry blind to groups, no lineage.}
\idea{Averages hide groups.}{Score each named group separately, with a confidence interval; a slice
  too small to bound its error reports \emph{insufficient evidence}.}
\idea{Fairness criteria are choices.}{With unequal base rates and imperfect prediction, calibration
  and equal error rates cannot both hold.}
\idea{Decisions select future labels.}{A denied applicant leaves no repayment record, so a gap can
  grow at every retraining.}
\idea{Responsibility costs the machine.}{Per-group labels, explanations, exploration traffic and audit logs
  cost compute; model cards and lineage make claims provable.}
\end{ideas}

\vspace{10pt}
\sheetsec{Fairness criteria, from a confusion matrix}
\begin{sheettable}{@{}l L L@{}}
\thead{Criterion} & \thead{Must match across groups} & \thead{Computed as}\\ \midrule
Demographic parity & approval rate & (TP + FP) $\div$ all\\
Equal opportunity & true positive rate (TPR) & TP $\div$ (TP + FN)\\
Equalized odds & TPR and false positive rate (FPR) & also FP $\div$ (FP + TN)\\
\end{sheettable}
\vspace{3pt}
{\fontsize{8.4}{10.4}\selectfont\color{muted}TP, FN: would repay, approved or rejected; FP, TN: would
default. Accuracy $=$ (TP + TN) $\div$ all. Gaps in percentage points (pts).\par}
\vspace{5pt}
On the fairness-accuracy \emph{Pareto frontier}, a smaller gap costs accuracy; a change that improves both
started below it. A lower cutoff raises TPR and FPR.\par
\end{leftcol}%
\begin{rightcol}
\sheetsec{Formulas you will use}
\fcard{Bias feedback model}%
  {\GroupDisparity(\FeedbackIntervalCount) \approx \GroupDisparity(0)\,\BiasFeedbackFactor^{\FeedbackIntervalCount}}%
  {Gap $\GroupDisparity$ after $\FeedbackIntervalCount$ retrainings (here one per quarter). Fit $\BiasFeedbackFactor =
   \GroupDisparity(1)/\GroupDisparity(0)$; for $\BiasFeedbackFactor > 1$ it first exceeds a gate once
   $\FeedbackIntervalCount > \ln(\text{gate}/\GroupDisparity(0)) / \ln\BiasFeedbackFactor$.}%
  {Sampling outcomes of marginal denials can pull $\BiasFeedbackFactor$ below one. Refit every audit.}%
  {\begin{tikzpicture}[sheet]
     \pgfmathsetmacro\sx{25}  \pgfmathsetmacro\sy{11}
     \foreach \q in {0,2,4,6,8}{\draw[gridline] ({20+\q*\sx/1},0) -- ({20+\q*\sx/1},74);
       \node[small, anchor=north] at ({20+\q*\sx/1},-3) {\q};}
     \node[small, anchor=north west] at (226,-3) {quarters};
     \draw[dim] (20,0) -- (220,0);
     \draw[accent, dashed, line width=0.6pt] (20,{\val{gate}*\sy}) -- (220,{\val{gate}*\sy});
     \node[note, anchor=south west] at (22,{\val{gate}*\sy+2}) {gate \val{gate} pts};
     \draw[accent, line width=1.1pt] (20,{\val{riseq0}*\sy})
       \foreach \q/\v in {1/\val{riseq1},2/\val{riseq2},3/\val{riseq3},4/\val{riseq4},5/\val{riseq5},6/\val{riseq6},7/\val{riseq7},8/\val{riseq8}}{ -- ({20+\q*\sx},{\v*\sy}) };
     \draw[muted, line width=0.8pt] (20,{\val{fallq0}*\sy})
       \foreach \q/\v in {1/\val{fallq1},2/\val{fallq2},3/\val{fallq3},4/\val{fallq4},5/\val{fallq5},6/\val{fallq6},7/\val{fallq7},8/\val{fallq8}}{ -- ({20+\q*\sx},{\v*\sy}) };
     \filldraw[fill=accent, draw=accent] ({20+\val{breachquarter}*\sx},{\val{breachexact}*\sy}) circle (2.2pt);
     \node[lab, anchor=south east] at ({17+\val{breachquarter}*\sx},{\val{breachexact}*\sy+3}) {breach: \val{breachgap} pts};
     \node[note, anchor=south] at ({20+2*\sx},{\val{riseq2}*\sy+3}) {$\BiasFeedbackFactor = \val{feedbackfactor}$};
     \node[note, anchor=south east, align=right] at (220,{\val{fallq8}*\sy+7}) {$\BiasFeedbackFactor = \val{sampledfactor}$:\\outcomes sampled};
     \node[small, anchor=east] at (17,{\val{launchgap}*\sy}) {\val{launchgap}};
     \node[small, anchor=south west] at (22,74) {TPR gap, pts};
   \end{tikzpicture}}
\fcard{Slice sample size}%
  {\text{margin} \approx \val{zscore}\sqrt{\val{worstvariance}/\SubgroupSampleSize}
   \qquad \RandomSampleCollectionSize = \SubgroupSampleSize \div \text{group share}}%
  {Worst-case \val{confidence}\% margin on a group's error rate from $\SubgroupSampleSize$ labeled examples,
   and what a random sample must collect to hold $\SubgroupSampleSize$ of them.}%
  {About \val{pergroup} examples give $\pm$\val{halfwidth} pt. A \val{groupshare}\% group then needs
   \val{randomtotal} at random (\val{randommultiple}$\times$), or \val{pergroup} recruited directly.}{}
\fcard{Lifecycle cost and carbon}%
  {\text{carbon} = \text{energy} \times \text{PUE} \times \text{carbon intensity}}%
  {Accelerator-hours $=$ requests $\times$ time each; cost $=$ hours $\times$ price. A \val{gpupower} kW
   accelerator on a \val{intensity} kg/kWh grid emits \val{kgpergpuhour} kg CO$_2$e per hour.}%
  {Here PUE $=$ 1: accelerator only, no host or manufacturing.}%
  {\begin{tikzpicture}[sheet]
     \pgfmathsetmacro\u{2.3}
     \pgfmathsetmacro\xa{\val{sharetraining}*\u} \pgfmathsetmacro\xb{\xa+\val{shareinference}*\u}
     \pgfmathsetmacro\xc{\xb+\val{shareoperations}*\u}
     \draw[hotbox] (0,0) rectangle (\xa,15);
     \draw[accbox] (\xa,0) rectangle (\xb,15); \node[lab] at ({(\xa+\xb)/2},7.5) {inference \val{shareinference}\%};
     \draw[nbox] (\xb,0) rectangle (\xc,15); \node[lab] at ({(\xb+\xc)/2},7.5) {operations \val{shareoperations}\%};
     \node[note, anchor=south west] at (0,17) {training \val{sharetraining}\%: \val{cycles} quarterly cycles};
     \node[note, anchor=south east] at (\xc,17) {DLRM, \val{years} years: \$\val{tcototal}M};
   \end{tikzpicture}}
\end{rightcol}

% ===================================================================== BACK: PRACTICE
\backside
\begin{answers}
\ans{Finish}{TPR gap \val{tpra} $-$ \val{newtpr} $=$ \val{newtprgap} pts: equal opportunity passes. FPR
  \val{newfp} $\div$ \val{unqualifiedb} $=$ \val{newfpr}\% vs.\ \val{fpra}\%, an \val{newfprgap}-pt gap: equalized odds
  fails. Approval \val{newapproved} $\div$ \val{totalb} $=$ \val{newapproval}\% vs.\ \val{approvala}\%: parity passes
  (\val{newparitygap} pts). The cost: more approved defaulters. Accuracy rose too: the old cutoff sat below the Pareto frontier.}
\ans{1}{$\BiasFeedbackFactor = \val{qoneaudit}/\val{qonelaunch} = \val{qonefactor}$; $\ln(\val{gate}/\val{qonelaunch})
  /\ln \val{qonefactor} = \val{qonelog}$: first exceeds it at quarter \val{qonequarters} (\val{qonebefore}, then
  \val{qonebreach} pts). At $\BiasFeedbackFactor = \val{sampledfactor}$ it
  falls to \val{qonesampled} pts and never breaches.}
\ans{2}{\val{qtworequests}M $\times$ \val{servicems} ms $=$ \val{qtwohoursday} accelerator-hours a day; $\times$ \val{days}
  days $\times$ \$\val{inferencerate} $=$ \$\val{qtwoinfcost}, \val{qtworatio}$\times$ training's \val{cycles} $\times$
  \$\val{cyclecost}. The serving cut saves \$\val{qtwoservingsave}, halving training \$\val{qtwotrainsave}. Carbon:
  \val{qtwohours} $\times$ \val{kgpergpuhour} kg $=$ \val{qtwoinfcarbon} t of emissions.}
\ans{3}{Algorithm: proxy and goal split, a gap of \val{qthreeclicklevel} $-$
  \val{qthreesatlevel} $=$ \val{qthreegap} of baseline. Matching inputs make drift unlikely; biased labels
  or a shift on the affected slice would point to Data. Measure satisfaction on a randomized holdout.}
\ans{4}{\val{pergroup} $\div$ \val{qfourshare}\% $=$ \val{qfourtotal} (\val{qfourmultiple}$\times$). With
  \val{qfourslice} examples the margin is \val{zscore}$\sqrt{\val{worstvariance}/\val{qfourslice}}$
  $=$ \val{qfourmargin} pts, wider than $\pm$\val{halfwidth}: insufficient evidence. Collect targeted labels.}
\end{answers}

\begin{leftcol}
\sheetsec{Worked example}
\begin{practicebox}\boxtitle{Does the credit service pass its release gate?}
Constructed counts; gate: TPR gap at most \val{gate} pts.\par\vspace{2pt}
\begin{sheettable}{@{}l C C C C@{}}
\thead{Group} & \thead{TP} & \thead{FN} & \thead{FP} & \thead{TN}\\ \midrule
A & \val{tpa} & \val{fna} & \val{fpa} & \val{tna}\\
B & \val{tpb} & \val{fnb} & \val{fpb} & \val{tnb}\\
\end{sheettable}
\par\vspace{6pt}
{\centering
\begin{tikzpicture}[sheet]
  \pgfmathsetmacro\u{1.9}
  \foreach \lab/\v/\y/\sty in {Group A TPR/\val{tpra}/18/accbox, Group B TPR/\val{tprb}/0/hotbox}{
    \node[lab, anchor=east] at (66,\y+6) {\lab};
    \draw[\sty] (70,\y) rectangle ({70+\v*\u},\y+12);
    \node[lab, anchor=west] at ({74+\v*\u},\y+6) {\v\%};}
  \draw[dim] ({70+\val{tprb}*\u},-3) -- ({70+\val{tprb}*\u},33);
  \draw[dim] ({70+\val{tpra}*\u},-3) -- ({70+\val{tpra}*\u},33);
  \node[note, anchor=south] at ({70+(\val{tpra}+\val{tprb})/2*\u},33) {gap \val{tprgap} pts};
\end{tikzpicture}\par}
\begin{steps}
\item Accuracy: (\val{correcta} + \val{correctb}) $\div$ \val{totalall} $=$ \res{\val{accuracyall}\%}.
\item Approval: A \val{approveda} $\div$ \val{totala} $=$ \val{approvala}\%, B \val{approvedb} $\div$ \val{totalb}
  $=$ \val{approvalb}\%: a \res{\val{paritygap}-pt parity gap}.
\item TPR: A \val{tpa} $\div$ \val{qualifieda} $=$ \val{tpra}\%, B \val{tpb} $\div$ \val{qualifiedb} $=$ \val{tprb}\%:
  \res{\val{tprgap} pts}, \val{gatemultiple}$\times$ the gate. FPR is \val{fpra}\% in both.
\end{steps}
\takeaway{Accuracy looks fine, yet the TPR gap fails the bound.}
\end{practicebox}

\vspace{12pt}
\sheetsec{Your turn: finish it}
\begin{practicebox}\boxtitle{Did lowering Group B's cutoff fix it?}
\assumed{} after the change, Group B has TP \val{newtp}, FN \val{newfn}, FP \val{newfp}, TN \val{newtn};
Group A is unchanged.\par\vspace{4pt}
\textbf{Given:} TPR $=$ \val{newtp} $\div$ \val{qualifiedb} $=$ \val{newtpr}\%.\par
\textbf{You:} find the TPR gap, Group B's FPR and approval rate. Which criteria now pass at a
\val{gate}-pt bound, and what did the change cost?
\end{practicebox}
\end{leftcol}%
\begin{rightcol}
\sheetsec{Common mistakes}
\mistake{``We dropped race and gender from the inputs, so it cannot discriminate.''}%
  {Proxies carry the signal: ZIP codes, college names, a cost target standing in for health need.
}{}
\mistake{``It passed the fairness gate at launch, so it is fair.''}%
  {Launched at \val{launchgap} pts under a \val{gate}-pt gate, it crosses at quarter
   \val{breachquarter} (\val{breachgap} pts) at $\BiasFeedbackFactor = \val{feedbackfactor}$, with no code change.}{}
\mistake{``We deleted the user's rows, so the erasure request is met.''}%
  {Weights can memorize a record; membership inference can reveal it was used. Lineage finds affected
   models to retrain or unlearn.}{}

\vspace{2pt}
\sheetsec{Check yourself}
\begin{qs}
\item \assumed{} a credit model launches with a \val{qonelaunch}-pt TPR gap; the next quarterly audit reads
  \val{qoneaudit} pts. When does it first exceed the \val{gate}-pt gate if the factor holds? If sampling refits the factor to
  \val{sampledfactor} from launch, what is the gap at that same quarter?
\item \assumed{} \val{qtworequests} million requests a day at \val{servicems} ms of accelerator time each, \$\val{inferencerate}
  per accelerator-hour, for \val{years} years; quarterly retraining at \$\val{cyclecost} a cycle. Which saves more:
  \val{qtwoservingcut}\% less time per request, or half the training? Total serving carbon?
\item \assumed{} a ranking change lifts clicks \val{qthreeclicks}\% while satisfaction falls \val{qthreesat}\%, and
  input distributions match training. Which D·A·M axis failed, and what evidence would point to another?
\item \assumed{} a group is \val{qfourshare}\% of users and needs \val{pergroup} labeled examples ($\pm$\val{halfwidth} pt).
  How many must a random sample collect? The release slice has \val{qfourslice}: what does the check report?
\end{qs}
\end{rightcol}
\end{document}
